\section{Previous Controller Iterations}
\label{sec:old-controller}

The current controller described in Section~\ref{sec:controllers} was preceded by two implementations, \texttt{OldController\_1} and \texttt{OldController\_2}, developed in sequence. Each iteration validated a specific part of the sense-think-act loop before the project committed to the modular architecture used today. This section records their design choices, what they demonstrated, and the concrete limitations that motivated the redesign, so that the evolution of the system remains traceable.

Beyond the control logic itself, each iteration also shaped the simulated environment it was tested against, from a small static arena to a warehouse world that eventually included obstacles capable of moving independently of the robot. As discussed below, that last change was not cosmetic: introducing moving obstacles altered what the perception and obstacle-avoidance code in \texttt{OldController\_2} actually had to guarantee, and is one of the clearest examples of the environment driving a design decision rather than the other way around.

\subsection{Initial Approach}
\label{subsec:old-controller-initial}

Development started on a Gctronic \textbf{e-puck} robot placed in a small test arena, and only later moved to a \textbf{Pioneer 3-DX} robot inside the full-scale warehouse world. This substitution was a change of simulated hardware and environment scale, independent of the two controller logic iterations described below: \texttt{OldController\_1} was written for the e-puck arena, while \texttt{OldController\_2} was written for the Pioneer 3-DX warehouse world that the project still uses today. The two robot models differ substantially in size, which is reflected directly in the kinematic constants each controller hardcodes: \texttt{OldController\_1} uses a wheel base of \texttt{0.052}\,m and a wheel radius of \texttt{0.02}\,m, while \texttt{OldController\_2} uses \texttt{0.33}\,m and \texttt{0.0975}\,m respectively, matching the Pioneer 3-DX manual. The change in scale is also visible in the goal coordinates each controller targets: \texttt{OldController\_1} drives toward a single hardcoded point near \texttt{(0.5, 0.5)}, consistent with a small test arena, whereas \texttt{OldController\_2} cycles through a list of warehouse-scale waypoints such as \texttt{(-29.3, 4)} and \texttt{(18.75, 2.25)}. Because the e-puck stage focused on validating the basic control and logging loop rather than warehouse-scale navigation, its lessons about obstacle-avoidance thresholds and acceleration behaviour do not translate directly to the Pioneer 3-DX scale, and had to be re-derived in \texttt{OldController\_2}. The warehouse world itself grew in the same step, from a small empty arena to a layout with fixed shelving and, as described below, moving obstacles that \texttt{OldController\_1} never had to contend with.

\texttt{OldController\_1} implements a single-file, purely reactive controller whose purpose was to establish the minimal working pipeline: read the lidar and GPS, decide a wheel velocity, and persist a run history through the \texttt{RobotLog} service, tagged with \texttt{controller\_version="old\_controller"} so its runs remain distinguishable in the logs from later iterations. The only quantity it carries between simulation steps is the current wheel velocity read back from the motors; the controller keeps no explicit heading estimate. At every timestep it splits the 360-point lidar image, read by \texttt{read\_lidar\_image}, into three equal sectors (left, center, right) and computes the minimum range in each, and it derives a bearing toward the fixed goal directly from the raw GPS position with \texttt{get\_angle}, rather than from any filtered state estimate. Sensor readings, wheel speeds and GPS coordinates are printed to the console roughly every half second, which was primarily a debugging aid for observing the controller live while the arena was small enough to watch directly. This console-first style of inspection was practical for a hand-sized arena and a single fixed goal, but it does not scale to a longer, multi-waypoint warehouse run in which relevant events are sparser in time and easy to miss in a scrolling terminal; this is part of the motivation for the more structured, periodically persisted logging that \texttt{OldController\_2} adds, described later in this section.

Motion is selected through a fixed cascade of distance thresholds on the three sector minima: a full stop-and-reverse for a detected dead end, a hard turn for a close frontal obstacle, a gradual deceleration for a more distant frontal obstacle, and a small corrective turn when only a side sector is triggered; in the absence of any obstacle the robot ramps up to a cruise speed. The \texttt{acc\_speed} helper applies linear acceleration and deceleration ramps toward the selected target speed so that velocity commands do not change discontinuously between steps, and \texttt{robot\_to\_wheel\_velocity} converts the resulting linear and angular velocity into individual wheel angular velocities using the wheel base and wheel radius, saturating both wheels together so that the requested curvature is preserved even when the maximum wheel speed would otherwise be exceeded. This approach proved effective for validating that lidar readings, GPS, motor commands and the logging pipeline (obstacle-state tracking, idle-time accounting, event logging) worked together correctly on a simple, single-goal task. It does not scale to the warehouse scenario, however: the goal is a single point hardcoded in the main loop, so the controller has no notion of a route or of multiple waypoints; obstacle avoidance is purely threshold-based on the current sector minima, with no memory of previously seen obstacles and no deliberate path planning, so the robot can only react once it enters an obstacle's immediate vicinity rather than anticipate it; and the distance thresholds (for example \texttt{0.35}\,m for a frontal obstacle) were tuned for the e-puck's scale and would need to be re-derived for a Pioneer 3-DX operating in a much larger environment, which is exactly what \texttt{OldController\_2} did.

\texttt{OldController\_2} restructures the control logic into a \texttt{RobotController\_v1} class built around explicit \texttt{State} and \texttt{Position} dataclasses, and introduces a list of sequential warehouse waypoints (\texttt{goal\_positions}) that the robot advances through and loops over once the final goal is reached. It reuses the same wheel-velocity conversion and saturation pattern as \texttt{OldController\_1} in \texttt{set\_robot\_velocity}, adapted to the Pioneer 3-DX's kinematic constants, and adds a camera device that is initialized and enabled but not yet consumed by the obstacle-avoidance logic, suggesting it was staged for a vision-based extension that was not completed in this iteration. Logging was extended accordingly: beyond the run history saved through \texttt{RobotLog}, the main loop periodically calls \texttt{save\_to\_database} roughly every five simulated seconds so that intermediate progress is persisted even if the simulation is interrupted.

The surrounding world was also extended with two \texttt{MovingWalls} instances, giving the controller dynamic obstacles to react to for the first time, in addition to the static warehouse layout. Each \texttt{MovingWalls} object locates a Webots node by its \texttt{DEF} name (\texttt{MOVING\_WALL\_1} and \texttt{MOVING\_WALL\_2}) through the controller robot's Supervisor interface, records that node's initial translation, and on every call to \texttt{move\_wall} rewrites the corresponding coordinate of the node's \texttt{translation} field as the initial position plus a sinusoidal offset with a configured amplitude and frequency along a chosen axis. In the main loop this method is invoked unconditionally at the very start of every simulation step, ahead of state estimation, planning and logging: the first wall oscillates along the \texttt{y} axis with an amplitude of \texttt{3}\,m, the second along the \texttt{x} axis with an amplitude of \texttt{5}\,m, both driven at the same angular frequency of \texttt{0.1}\,rad/s. Because \texttt{getFromDef} and direct field mutation are Supervisor-only operations in Webots, running this code requires the \texttt{OldController\_2} process itself to hold Supervisor privileges, so a single script ends up simultaneously responsible for driving the vehicle and for choreographing part of the environment it drives through, a coupling that \texttt{OldController\_1} never needed since its arena was entirely static.

Critically, \texttt{RobotController\_v1} has no dedicated code path for the moving walls: from the lidar's perspective a moving wall is indistinguishable from any other obstacle, and its motion only shows up indirectly, as occupied sector ranges that appear, disappear or shift between consecutive scans faster than a purely static obstacle's would. The two \texttt{MovingWalls} are therefore the concrete scenario in this iteration that exercises the sector-clustering logic's persistent-identity bookkeeping, described in the obstacle-avoidance discussion below, even though that logic itself does not treat static and dynamic obstacles any differently. In other words, the moving walls changed the requirements placed on the perception and avoidance code without requiring a single line written specifically for them: the obstacle-avoidance layer had to be correct for a world where sector boundaries can move on their own, and the walls were the mechanism by which that requirement was actually tested during development.

State estimation in \texttt{OldController\_2} combines wheel odometry, integrated through \texttt{update\_wheel\_odometry} and \texttt{update\_with\_odometry}, with GPS correction in \texttt{fuse\_with\_gps}. In practice this fusion is configured with a GPS weight of \texttt{1.0}, so the corrected position is always the raw GPS reading and the odometry contribution is overwritten at every step rather than blended with it; the odometry-based prediction is computed but not actually used to stabilize the estimate. Heading and distance errors to the current goal are converted into linear and angular velocity commands by \texttt{calculate\_velocity}, a proportional controller with gains \texttt{K\_DISTANCE} and \texttt{K\_HEADING} and a \texttt{cos(alpha)} term that slows the robot down when the heading error is large, with the resulting linear and angular velocities clamped to configured maxima and a small dead-zone applied to the angular command to avoid steady-state oscillation. Because this proportional scheme reacts only to the instantaneous heading and distance error, it carries no memory of previous corrections and no explicit model of the vehicle's dynamics, which is consistent with the controller's overall character as a direct, low-latency reaction to the current state estimate rather than a trajectory-tracking scheme.

Obstacle avoidance is the most developed part of this iteration. The lidar field of view is divided into 32 sectors, each classified as free or occupied, and consecutive occupied or free sectors are grouped into obstacle and free-space \emph{clusters} that are assigned persistent identifiers across timesteps (\texttt{used\_obstacle\_ids}, \texttt{used\_space\_ids}, \texttt{previous\_scan}), so that the same physical obstacle or gap keeps the same identity from one lidar scan to the next even as its exact sector boundaries shift. The two \texttt{MovingWalls} described above are the clearest source of this kind of shift: a purely static obstacle's sector boundaries change only as the robot itself moves, whereas a moving wall's boundaries also change on their own, independently of the vehicle's trajectory, so the persistent-identity mechanism has to hold up under a case the static warehouse layout alone would not exercise. The controller then locks onto a specific free-space cluster (\texttt{locked\_space}) while an obstacle blocks the direct path to the goal, which prevents it from oscillating between two competing free sectors on successive frames, and releases the lock once the obstacle is no longer relevant. When no free sector remains, the controller enters an escape mode that rotates in place until a free sector reappears. This sector-clustering strategy proved robust enough that, as discussed below, it was carried forward almost unchanged into the current controller.

The implementation also exposes concrete limitations. The persistent-identity logic that tracks clusters across frames is intricate and stateful, and the code raises an exception if a sector is left without an assigned identifier, which is a fragile failure mode for a long-running simulation. The file still contains earlier, commented-out obstacle-avoidance attempts (\texttt{old\_obstacle\_avoidance}, \texttt{min\_distances}), which indicates that the sector-based strategy was reached iteratively rather than designed upfront. Goal waypoints and robot-specific parameters remain hardcoded in the controller script instead of being externalized to a configuration file. The Supervisor privileges required to animate the \texttt{MovingWalls} are a further limitation: granting a per-robot controller direct write access to arbitrary scene nodes is difficult to justify once several AGVs share the same warehouse, since any one controller process could in principle reposition world geometry that every other vehicle relies on for localization and obstacle avoidance, and \texttt{OldController\_2} does not scope or guard that access in any way.

\subsection{Lessons Learned}
\label{subsec:old-controller-lessons}

Comparing the two iterations with the current controller (Section~\ref{sec:controllers}) shows a clear line of continuity rather than a full rewrite. The sector-clustering obstacle-avoidance algorithm introduced in \texttt{OldController\_2} is, in its essential structure, the same algorithm implemented today by the \texttt{SectorPlanner} low-level planner: the same notion of dividing the lidar field of view into sectors, assigning persistent obstacle and free-space identifiers across frames, locking onto a chosen free sector, and rotating to escape a dead end. This is direct evidence that the local, reactive avoidance layer designed and tuned in \texttt{OldController\_2} was robust enough to be promoted into production, once decomposed behind a dedicated planner interface and driven by configuration objects instead of hardcoded constants. The \texttt{MovingWalls} introduced in \texttt{OldController\_2} were, in this sense, an early stress test for that design: because they could enter or leave a sector independently of the robot's own trajectory, they exercised the persistent-identity bookkeeping in a way a purely static layout would not, and the fact that the same clustering approach was carried forward into the \texttt{SectorPlanner} suggests it held up under that test rather than being reworked to cope with dynamic obstacles later. The lower-level wheel-velocity conversion pattern first written in \texttt{OldController\_1} and reused in \texttt{OldController\_2} survived the transition as well, which suggests that the basic differential-drive actuation model was correct from the earliest iteration and needed no substantial rework, only re-parameterization for each robot's physical dimensions.

Other aspects of the design changed more substantially. Route information moved out of the controller code entirely: instead of the hardcoded \texttt{goal\_positions} list, waypoints and per-robot routes are now declared in \texttt{goals.config.json} as named locations, which decouples warehouse layout changes from controller logic and lets new pickup or dropoff points be added without touching Python code. The monolithic \texttt{RobotController\_v1} class was split into dedicated perception, localization, planning and hardware modules, and a high-level planner was introduced above the reactive sector layer so the vehicle now follows a globally planned route rather than only locally reactive corrections. Hardware access was also wrapped behind a dedicated hardware-interface layer instead of calling \texttt{robot.getDevice} directly from control logic, improving testability and making the sensor and motor interfaces explicit; the camera device that \texttt{OldController\_2} enabled but never used is, in this later architecture, finally exercised by a dedicated perception module.

The GPS/odometry fusion limitation identified in \texttt{OldController\_2} was also addressed directly rather than carried forward silently: the current localization module generalizes the fixed \texttt{gps\_weight = 1.0} into a configurable weighted fusion, adds a disagreement threshold that snaps the state estimate to GPS only when odometry has drifted too far, and additionally fuses a heading estimate from the inertial unit, which the old controllers did not use at all. Taken together, these observations show that the old controllers were most useful not as discarded prototypes but as a way to isolate which parts of the design were already sound and which parts still needed rework: the sector-based reactive avoidance layer and the differential-drive actuation model proved robust enough to keep, while state estimation, goal configuration and, most importantly, coordination between multiple AGVs sharing the same warehouse still required further work before the system could operate reliably at full scale.
